\section{Metodología}
    El desarrollo se implementó en dos scripts independientes que se ejecutan con orden
    incremental, puesto que cada paso depende del anterior, iniciando con la preparación
    del entorno, siguiendo con la generación del conjunto de datos, la construcción de la
    distribución de frecuencias por la regla de Sturges, el cálculo de las medidas de
    tendencia central y de dispersión y, por último, la producción de los gráficos
    básicos y su verificación cruzada en R. Ambos scripts se reproducen completos en los
    Anexos \ref{anexo_python} y \ref{anexo_r}, de modo que esta sección describe las
    decisiones de implementación y el lector puede contrastarlas allí línea por línea.

    \subsection{Herramientas y entorno de trabajo}
        Las alternativas disponibles para un laboratorio de estadística descriptiva 
        se agrupan en tres familias, las hojas de cálculo, que ofrecen funciones y 
        gráficos inmediatos pero mezclan datos, fórmulas y resultados en un mismo 
        lienzo y dificultan auditar los cálculos y reproducirlos sobre datos nuevos
        \cite{toro2022ciencia}, \cite{timbers2024datascience},
        \cite{castro2026cienciadatos}, los entornos de programación científica, que
        expresan cada cálculo como código versionable capaz de regenerar exactamente 
        los mismos resultados y figuras \cite{castro2026visualizacion} y las librerías
        declarativas junto con las plataformas de inteligencia de negocios, que 
        producen gráficos con poco código a cambio de un menor control sobre los 
        elementos que los principios de diseño exigen ajustar \cite{few2012show}.

        \vspace{10pt}

        Se elige la segunda familia con \textit{pandas} y \textit{Matplotlib} como
        herramienta principal por dos razones, la primera es que el cálculo y su
        representación conviven en el mismo script, pues la tabla de frecuencias 
        se obtiene con \texttt{numpy.histogram} \cite{harris2020numpy} y las medidas 
        por grupo con \textit{pandas} \cite{mckinney2010data} y esas mismas estructuras
        alimentan directamente las figuras, eliminando transcripciones manuales que son
        fuente de error, la segunda es que \textit{Matplotlib} otorga el control fino 
        que exige la aplicación de los principios de diseño \cite{hunter2007matplotlib}, 
        ya que permite fijar las clases de Sturges como bordes del histograma, superponer
        las líneas de media, mediana y moda, forzar el eje de frecuencias desde cero o
        situar la cuadrícula detrás de los datos \cite{few2012show}, 
        \cite{castro2026visualizacion}. Como herramienta de contraste se emplea R base, 
        que implementa los mismos cálculos con funciones nativas como \texttt{mean}, 
        \texttt{median}, \texttt{var}, \texttt{sd} e \texttt{hist}, sin dependencias 
        externas \cite{castro2026visualizacion}, \cite{rcore2024r}. La comparación es 
        directa porque \texttt{var()} y \texttt{sd()} de R son muestrales, con divisor 
        $n - 1$, exactamente como el parámetro \texttt{ddof=1} de \textit{pandas}, 
        de modo que los resultados deben coincidir dígito a dígito.

    \subsection{Conjunto de datos}
        Se emplea un conjunto de datos simulado que corresponde al consumo energético
        mensual de 120 clientes de una empresa distribuidora, generado con semilla fija
        (42) para garantizar que las observaciones sean reproducibles y puedan
        regenerarse de manera idéntica en cada ejecución. Los sectores se asignan con
        probabilidades de 0,50 para Residencial, 0,30 para Comercial y 0,20 para
        Industrial, el consumo proviene de distribuciones normales con parámetros
        propios de cada sector y el costo resulta de una tarifa diferenciada con un
        ruido multiplicativo del 4\%. Las variables del dataset se muestran a 
        continuación:

        \begin{table}[H]
            \centering
            \caption{Variables del conjunto de datos}
            \begin{center}
                \begin{tabular}{|l|l|>{\raggedright\arraybackslash}p{3.1cm}|}
                    \hline
                    \textbf{Variable} & \textbf{Tipo} & \textbf{Descripción} \\
                    \hline
                    cliente\_id & Nominal & Identificador único del cliente
                    (CL-001 a CL-120) \\
                    \hline
                    sector & Nominal & Residencial, Comercial o Industrial \\
                    \hline
                    consumo\_kwh & Cuantitativa continua & Consumo mensual en
                    kilovatios-hora \\
                    \hline
                    costo\_miles\_cop & Cuantitativa continua & Facturación mensual en
                    miles de pesos colombianos \\
                    \hline
                \end{tabular}
                \label{dataset_variables}
            \end{center}
        \end{table}

        El carácter simulado del dataset constituye una ventaja metodológica para este
        laboratorio, puesto que al conocerse de antemano que existen tres poblaciones
        con medias separadas, es posible evaluar si cada estadístico y cada gráfico
        revela u oculta esa estructura, explicando con precisión por qué las medidas
        globales difieren tanto de las sectoriales.

    \subsection{Flujo de trabajo y reproducibilidad}
        El flujo inicia en Python (Anexo \ref{anexo_python}), donde \textit{statistics.py} 
        genera el dataset, calcula los estadísticos y exporta los resultados y las 
        figuras en formatos CSV y PNG, posteriormente, R (Anexo \ref{anexo_r}) utiliza 
        el mismo conjunto de datos para recalcular los estadísticos de forma independiente y
        replicar las figuras mediante su graficación base, permitiendo contrastar ambos 
        resultados.

    \subsection{Construcción de la distribución de frecuencias}
        La tabla de frecuencias se construye aplicando la regla de Sturges 
        (ecuación \ref{sturges}) al consumo global para determinar el número de clases, 
        los intervalos, de igual amplitud, se generan con \texttt{numpy.linspace}, 
        mientras que \texttt{numpy.histogram} calcula las frecuencias absolutas para 
        luego obtener las frecuencias acumuladas y relativas y la marca de clase de 
        cada intervalo.

    \subsection{Cálculo de las medidas de tendencia central}
        La media y la mediana se calculan directamente con \textit{pandas}, mientras 
        que la moda, al tratarse de una variable continua, se estima mediante la 
        función \texttt{interpolated\_mode}, adicionalmente, se identifica la clase 
        modal con \texttt{argmax} y aplica la ecuación (\ref{moda_interpolada}), 
        considerando el caso $d_{1}+d_{2}=0$, en el que se utiliza el punto medio 
        de la clase. El procedimiento se repite para cada sector recalculando sus 
        intervalos según la regla de Sturges para adaptarlos al tamaño de cada 
        muestra.

    \subsection{Cálculo de las medidas de dispersión}
        Las medidas de dispersión se calculan con las ecuaciones (\ref{rango}) a
        (\ref{iqr}), fijando de manera explícita \texttt{ddof=1} en la varianza y en la
        desviación estándar para obtener los estimadores muestrales, que son justamente
        los que R implementa por defecto en \texttt{var()} y \texttt{sd()}, condición que
        hace posible la verificación cruzada dígito a dígito.

    \subsection{Verificación cruzada en R}
        El script de R no reutiliza ningún resultado de Python, únicamente lee el 
        archivo CSV con las 120 observaciones y vuelve a calcular desde cero la tabla 
        de frecuencias y todos los estadísticos, de manera que cualquier discrepancia
        revelaría un error de implementación en alguno de los dos entornos; la 
        distribución de frecuencias en R reproduce el mismo procedimiento, generando 
        los bordes con \texttt{seq} y contando con \texttt{hist} en modo silencioso, 
        con el argumento \texttt{include.lowest} activo para que la observación mínima 
        quede dentro de la primera clase, tal como lo hace \texttt{numpy.histogram}, 
        por último, la función \texttt{summary\_stats} condensa en un solo vector las
        ocho medidas de las Tablas \ref{central_tendency} y \ref{dispersion} y se 
        aplica a cada sector y al conjunto global.